% Partie : théorie de la démonstration
% Chapitre : déduction naturelle
% Section : règles et preuves

\documentclass[../../../include/open-logic-section]{subfiles}

\begin{document}

\olfileid[fr]{pt}{nat}{rp}

\olsection{Règles et \usetoken{p}{proof}}

Commençons par quelques définitions et précisons la terminologie.

Considérons les deux règles de \Log{N1c} et de~\Log{N1i} qui
font intervenir $!A \lif !B$.
\[
\bottomAlignProof
\AxiomC{$\Discharge{!A}{x}$}
\shortDeduce
\DeduceC{$!B$}
\DischargeRule{\Intro{\lif}}{x}
\UnaryInfC{$!A \lif !B$}
\DisplayProof
\qquad
\bottomAlignProof
\AxiomC{$!A \lif !B$}
\AxiomC{$!A$}
\RightLabel{\Elim{\lif}}
\BinaryInfC{$!B$}
\DisplayProof
\]
La règle \Elim{\lif} est simple. Les !!{formula}s situées
au-dessus du trait sont les \emph{prémisses}. Celle de gauche,
$!A \lif !B$, est la !!{formula} \emph{principale}. Chaque
règle d'!!{elimination} possède une telle prémisse, toujours
à l'extrême gauche : c'est la prémisse \emph{majeure} de
l'inférence. Les autres prémisses, lorsqu'il y en a, comme ici,
sont dites \emph{mineures}. Les prémisses des règles
d'!!{introduction} sont également appelées prémisses majeures.
La !!{formula} située sous le trait est la \emph{conclusion}.

La règle \Intro{\lif} ne possède qu'une prémisse ; sa conclusion
$!A \lif !B$ est la formule principale. Il en va ainsi de
toutes les règles d'!!{introduction} : la conclusion est la
formule principale et son opérateur principal est celui que
l'on « introduit ».

La notation $\Discharge{!A}{x}$ s'interprète comme suit.
Dans !!a{proof}, l'inférence \Intro{\lif} s'applique à une
prémisse~$!B$. La sous-!!{proof} qui se termine par~$!B$ comporte
des !!{formula}s initiales, appelées \emph{hypothèses}.
!!^a{proof} de~$!B$ à partir d'une ou de plusieurs hypothèses
de la forme~$!A$ établit que $!A$ entraîne~$!B$, sous les
autres hypothèses éventuellement encore ouvertes.
L'application de \Intro{\lif} donne alors~$!A \lif !B$.
Cette conclusion ne dépend plus des occurrences de
l'hypothèse~$!A$ que l'on a choisies de décharger. Ces
occurrences sont dites \emph{déchargées} ou \emph{fermées}
dans la !!{proof} incluant l'inférence~\Intro{\lif}.
La règle précise quelles hypothèses peuvent être déchargées.
Pour \Intro{\lif}, avec conclusion $!A \lif !B$, ce sont les
hypothèses de la forme~$!A$, c'est-à-dire celles qui coïncident
avec l'antécédent de l'implication. L'étiquette de
décharge~$x$ permet d'indiquer quelles occurrences sont
déchargées par quelle inférence. Une règle qui décharge des
hypothèses n'est pas tenue de décharger \emph{toutes} les
occurrences de la forme requise ; elle peut même n'en
décharger aucune. Par exemple, la !!{proof} suivante est
correcte, bien que la première inférence \Intro{\lif} ne
décharge aucune hypothèse :
\[
\AxiomC{$\Discharge{!A}{y}$}
\RightLabel{\Intro{\lif}}
\UnaryInfC{$!B \lif !A$}
\DischargeRule{\Intro{\lif}}{y}
\UnaryInfC{$!A \lif (!B \lif !A)$}
\DisplayProof
\]
Lorsqu'une règle ne décharge aucune hypothèse, on peut omettre
l'étiquette de décharge, comme dans cette première inférence.

Dans la !!{proof}
\[
\AxiomC{$\Discharge{!A}{x}$}
\AxiomC{$\Discharge{!A}{y}$}
\RightLabel{\Intro{\land}}
\BinaryInfC{$!A \land !A$}
\RightLabel{\Elim{\land}}
\UnaryInfC{$!A$}
\DischargeRule{\Intro{\lif}}{x}
\UnaryInfC{$!A \lif !A$}
\DischargeRule{\Intro{\lif}}{y}
\UnaryInfC{$!A \lif (!A \lif !A)$}
\DisplayProof
\]
la première inférence \Intro{\lif} décharge l'hypothèse de
gauche~$!A^x$, et la seconde celle de droite~$!A^y$.
Toutes les hypothèses portant l'étiquette~$x$ sont déchargées
par l'inférence \Intro{\lif} étiquetée~$x$ ; celles
d'étiquette~$y$ le sont par l'inférence étiquetée~$y$.
Cela exige que toutes les hypothèses de même étiquette
soient aussi une même !!{formula}.

Les règles de quantification sont un peu plus complexes.
Par exemple, les règles de \Log{N1i} pour $\lexists$ sont :
\[
\bottomAlignProof
\AxiomC{$!A(t)$}
\RightLabel{\Intro{\lexists}}
\UnaryInfC{$\lexists[x][!A(x)]$}
\DisplayProof
\qquad
\bottomAlignProof
\AxiomC{$\lexists[x][!A(x)]$} 
\AxiomC{$\Discharge{!A(c)}{x}$}
\shortDeduce
\DeduceC{$!B$}
\DischargeRule{\Elim{\lexists}}{x}
\BinaryInfC{$!B$} \DisplayProof 
\bottomAlignProof
\]
Dans \Intro{\lexists}, la prémisse est~$!A(t)$, où $t$ est
un terme clos, c'est-à-dire sans variable. L'inférence est
correcte --- elle est une instance du schéma
\Intro{\lexists} --- s'il existe une !!{formula}~$!A$,
une variable~$x$ et un terme clos~$t$ tels que la conclusion
soit $\lexists[x][!A]$ et la prémisse~$\Subst{!A}{t}{x}$.
La règle \Elim{\lexists} permet de décharger des hypothèses
de la forme~$!A(c)$ : à chaque inférence correspond une
!!{constant}~$c$, et les hypothèses déchargées sont de la
forme $\Subst{!A}{c}{x}$. Elles doivent toutes employer la
même constante~$c$. Celle-ci est appelée \emph{constante
propre}, ou \emph{eigenvariable}, de l'inférence
\Elim{\lexists}. L'inférence n'est correcte que si $c$
est absente de toute hypothèse ouverte autre que les
occurrences désignées de~$!A(c)$ qui sont déchargées,
ainsi que de~$!B$ et de~$!A$. C'est la \emph{condition de
la constante propre}.

Les !!{proof}s de \Log{N1c} et de~\Log{N1i} sont des arbres
finis de formules, construits inductivement à partir des
hypothèses au moyen des règles d'inférence. Chaque feuille
est une hypothèse, éventuellement munie d'une étiquette de
décharge. Toute autre !!{formula} s'obtient par une règle du
système à partir des !!{formula}s immédiatement au-dessus.
Une occurrence d'hypothèse est dite \emph{ouverte} à une
inférence si aucune inférence située sur la branche qui
relie cette occurrence à l'inférence considérée ne l'a
déchargée.

\begin{defn}\ollabel{defn:proof}
Les \emph{!!{proof}s} de \Log{N1c} et de~\Log{N1i} sont
les arbres finis de !!{formula}s définis inductivement
comme suit.
\begin{enumerate}
  \item\ollabel{defn:prf:assumption} Toute !!{formula},
  munie d'une étiquette de décharge telle que $x$ ou $y$,
  constitue à elle seule !!a{proof}. Dans !!a{proof}
  réduite à une !!{formula}, celle-ci est une hypothèse ouverte.
  \item\ollabel{defn:prf:one-prem} Soit $R$ une règle à une,
  deux ou trois prémisses $!A_1$, $!A_2$, $!A_3$, de
  conclusion~$!A$. Soient $\delta_1$, $\delta_2$,
  $\delta_3$ des !!{proof}s se terminant respectivement
  par $!A_1$, $!A_2$, $!A_3$. Alors les arbres
\[
  \AxiomC{}
  \RightLabel{$\delta_1$}
  \DeduceC{$!A_1$}
  \RightLabel{$R$}
  \UnaryInfC{$!A$}
  \DisplayProof
\qquad
  \AxiomC{}
  \RightLabel{$\delta_1$}
  \DeduceC{$!A_1$}
  \AxiomC{}
  \RightLabel{$\delta_2$}
  \DeduceC{$!A_2$}
  \RightLabel{$R$}
  \BinaryInfC{$!A$}
  \DisplayProof
\qquad
  \AxiomC{}
  \RightLabel{$\delta_1$}
  \DeduceC{$!A_1$}
  \AxiomC{}
  \RightLabel{$\delta_2$}
  \DeduceC{$!A_2$}
  \AxiomC{}
  \RightLabel{$\delta_3$}
  \DeduceC{$!A_3$}
  \RightLabel{$R$}
  \TrinaryInfC{$!A$}
  \DisplayProof
  \]
  sont respectivement des !!{proof}s, pourvu que les
  étiquettes d'hypothèses soient compatibles --- deux
  !!{formula}s différentes ne doivent pas porter la même
  étiquette de décharge --- et que les inférences soient
  des applications correctes de~$R$.
  \item Les occurrences initiales de !!{formula}s dans
  la nouvelle !!{proof} sont des \emph{hypothèses ouvertes}
  si elles sont ouvertes dans $\delta_1$, $\delta_2$
  ou~$\delta_3$, sauf celles que la règle décharge.
  Lorsque la règle porte l'étiquette~$x$ (ou $x\, y$),
  \Intro{\lif}, \Intro{\lnot} et \FalseCl déchargent
  les hypothèses ouvertes de~$\delta_1$ étiquetées~$x$ ;
  \Elim{\lexists} décharge celles de~$\delta_2$
  étiquetées~$x$ ; \Elim{\lor} décharge celles
  de~$\delta_2$ étiquetées~$x$ et celles de~$\delta_3$
  étiquetées~$y$.
\end{enumerate}
\end{defn}

Les !!{proof}s employées dans la construction inductive
d'!!a{proof} sont ses \emph{sous-!!{proof}s}. Celles qui se
terminent par les prémisses d'une inférence sont les
\emph{sous-!!{proof}s immédiates} de cette inférence.
Les règles de \Log{N1c} et de \Log{N1i} figurent dans le
\olref{tab:N1}.

\subfile{rules-N1}

\begin{defn}
La \emph{!!{height}}~$\pheight{\delta}$ d'!!a{proof}~$\delta$
est définie inductivement comme suit.
\begin{enumerate}
  \item Si $\delta$ se réduit à une hypothèse,
  $\pheight{\delta} = 0$.
  \item Si $\delta$ se termine par une inférence à une
  prémisse, de sous-!!{proof} immédiate~$\delta_1$, alors
  $\pheight{\delta} = \pheight{\delta_1} + 1$.
  \item Si $\delta$ se termine par une inférence à deux
  prémisses, de sous-!!{proof}s immédiates~$\delta_1$
  et~$\delta_2$, alors $\pheight{\delta} =
  \max(\pheight{\delta_1}, \pheight{\delta_2}) + 1$.
  \item Si $\delta$ se termine par \Elim{\lor}, de
  sous-!!{proof}s immédiates~$\delta_1$, $\delta_2$
  et~$\delta_3$, alors $\pheight{\delta} =
  \max(\pheight{\delta_1}, \pheight{\delta_2}, \pheight{\delta_3}) + 1$.
\end{enumerate}
Pour un séquent~$S$ exprimant la conclusion et les hypothèses
ouvertes d'!!a{proof}, on écrit $\Proves[n] S$ lorsqu'il
existe une telle !!{proof} de hauteur $\le n$.
\end{defn}

\begin{ex}
Dans \Log{N1i}, toute !!{formula} constitue !!a{proof} à
partir d'elle-même prise comme hypothèse ouverte. Ainsi,
\[
(!A \land !B)^x
\]
est !!a{proof}~$\delta_1$, de hauteur~$0$.
On en tire deux !!{proof}s~$\delta_2$ et~$\delta_3$ en
appliquant respectivement les deux versions de~\Elim{\land} :
\[
\AxiomC{$(!A \land !B)^x$}
\RightLabel{\Elim{\land}[2]}
\UnaryInfC{$!B$}
\DisplayProof
\qquad
\AxiomC{$(!A \land !B)^x$}
\RightLabel{\Elim{\land}[1]}
\UnaryInfC{$!A$}
\DisplayProof
\]
La sous-!!{proof} immédiate de la dernière inférence
de~$\delta_2$ et de~$\delta_3$ est la même : $!A \land !B$.
Les deux !!{proof}s ont cette occurrence de $!A\land !B$
pour hypothèse ouverte. On a
$\pheight{\delta_2} = \pheight{\delta_3} = 1$.

La règle \Intro{\land} permet, à partir de deux prémisses
$!B$ et $!A$, de conclure~$!B \land !A$. L'arbre suivant
est donc !!a{proof}~$\delta_4$ :
\[
\AxiomC{$(!A \land !B)^x$}
\RightLabel{\Elim{\land}[2]}
\UnaryInfC{$!B$}
\LeftSubproofLabel{$\delta_2$}
\AxiomC{$(!A \land !B)^x$}
\RightLabel{\Elim{\land}[1]}
\UnaryInfC{$!A$}
\RightSubproofLabel{$\delta_3$}
\RightLabel{\Intro{\land}}
\BinaryInfC{$!B \land !A$}
\DisplayProof
\]
Sa hauteur est $\pheight{\delta_4} = \max(\pheight{\delta_2},
\pheight{\delta_3})+1 = \max(1,1) + 1 = 2$.
Ses deux occurrences de $!A \land !B$ sont des
hypothèses ouvertes.

Enfin, on applique \Intro{\lif} pour obtenir $\delta_5$ :
\[
\AxiomC{$\Discharge{!A \land !B}{x}$}
\RightLabel{\Elim{\land}[2]}
\UnaryInfC{$!B$}
\LeftSubproofLabel{$\delta_2$}
\AxiomC{$\Discharge{!A \land !B}{x}$}
\RightLabel{\Elim{\land}[1]}
\UnaryInfC{$!A$}
\RightSubproofLabel{$\delta_3$}
\RightLabel{\Intro{\land}}
\BinaryInfC{$!B \land !A$}
\DischargeRule{\Intro{\lif}}{x}
\UnaryInfC{$(!A \land !B) \lif (!B \land !A)$}
\DisplayProof
\]
Sa hauteur est $\pheight{\delta_4} + 1 = 3$.
L'inférence \Intro{\lif} décharge toutes les occurrences
d'hypothèses de la forme $!A \land !B$ étiquetées~$x$,
c'est-à-dire les deux hypothèses jusque-là ouvertes de
sa sous-preuve immédiate. Ce sont les seules hypothèses
de $\delta_5$ ; celle-ci n'a donc plus d'hypothèse ouverte.
Elle est une !!{proof} \emph{de} sa !!{formula} finale
$(!A \land !B) \lif (!B \land !A)$.
\end{ex}















\end{document}
